\subsection{Multi-Domain Fusion Benefits}
\label{sec:discussion-fusion}

The ablation results from Section~\ref{sec:results-ablation} demonstrate that multi-domain fusion is the single most impactful architectural choice in PMDF-Net. Removing the time-frequency branch and fusion module (PMDF-Net w/o multi-domain) degrades SNR by $2.1$ dB relative to the full model, while time-domain-only operation loses $2.6$ dB and time-frequency-only operation loses $3.8$ dB. These gaps are not marginal; in LU signal processing, a $2.6$ dB SNR improvement translates directly to detectable defect sizes that would otherwise fall below the instrument resolution floor.

The underlying reason is complementary representation. The time domain captures temporal ordering and transient waveform morphology, including abrupt phase shifts and localized pulse distortions that characterize narrowband defects. The frequency domain encodes spectral content --- energy distribution across harmonic bands, modal frequencies, and resonant signatures that are invariant to the precise timing of transient events. A defect that displaces a reflector's arrival time by a fraction of a microsecond produces a measurable time-domain phase anomaly, while the same defect may simultaneously shift spectral energy among resonant modes by an amount detectable in the frequency domain. Neither view alone preserves sufficient information to distinguish defect-related changes from noise-induced artefacts.

This complementarity is especially critical for defect detection. Many physically meaningful defect signatures do not dominate a single domain; they are split across both. A delamination that reduces interface stiffness subtly broadens a resonant peak without significantly altering the arrival time of the primary echo. Conversely, a fatigue crack that introduces a secondary micro-echo perturbs the time-domain waveform in a way that leaves spectral features largely intact. The fusion module in PMDF-Net learns to weight and align these domain-specific features such that weak signals in either domain are amplified by strong signals in the other, enabling detection decisions that would fail under either single-domain configuration. The $2.6$ dB gain from fusion is therefore not merely an SNR metric; it is a practical improvement in the lower bound of detectable flaw severity.